No scalar objective was optimised; instead the non-dominated sets of six objective pairs were extracted from the full database (Fig.~\ref{fig:pareto}, all structures; the porous-only fronts are in \texttt{figures/\allowbreak campaign/\allowbreak pareto\_fronts\_porous.pdf}). The pristine sheets dominate every front that does not reward porosity, so the informative fronts are those restricted to porous designs. On the strength--failure-strain front the load-parallel slit designs (25.4--25.6\,N/m) occupy the strong end, the 45$^\circ$ slit lattice the ductile end (0.32 failure strain through geometric re-stiffening), and the low-porosity mesh in between; on strength--work only the two parallel-slit designs are non-dominated; on stiffness--failure strain the front adds the $0/90^\circ$ strut lattice and the hierarchy $\times$ anisotropy mesh (17.5\,N/m, 0.26 failure strain); on work--load retention the coarse meshes, the three-level mesh and the slits share the front; on specific work per atom versus porosity the sparse strut lattices ($\phi=0.45$) are non-dominated despite low strength. The fronts are structurally diverse (slits, coarse meshes, nested anisotropic meshes, strut lattices, a shear-like slit lattice) and no family dominates more than two objective pairs; the trade-off most relevant to design is strength versus progressivity: the strongest porous designs (straight equivalent strips) fail in one avalanche, and the only designs that combine a strength above 16\,N/m with progressive failure are those with many straight, load-parallel ligaments of moderate dispersion (load-parallel-vein mesh, hierarchy $\times$ anisotropy mesh, 32\,\AA{} round-pore mesh).
